\FormSection{1. Thông tin chung:}
\FormLine{Mã số đề tài: \ProjectCode}
\FormLine{Tên đề tài: \ProjectTitle}
\FormLine{Sinh viên chủ nhiệm đề tài: \PrincipalName}
\FormLine{Lớp: \PrincipalClass{} -- Khóa \PrincipalCohort; Trường/Khoa/Viện: \PrincipalUnit.}
\FormIndented{Năm thứ: \PrincipalYear\qquad Số năm đào tạo: \TrainingYears}
\FormLine{Người hướng dẫn: \SupervisorName}
\FormSection{2. Mục tiêu đề tài:}
Thiết kế và hiện thực kiến trúc đa thuê bao cho Kanban; phân tích định danh, cô lập dữ liệu, quyền, nhất quán và quản lý tài nguyên. Xây dựng chuỗi truy vết từ bài toán tới bằng chứng và phương pháp đánh giá.
\FormSection{3. Tính mới và sáng tạo:}
Đóng góp sơ bộ là kiến trúc tham chiếu và ma trận kiểm chứng trên cùng nghiệp vụ. Tính mới học thuật và hiệu quả so sánh cần được xác lập qua tổng quan, protocol và dữ liệu; chưa tuyên bố đã xác nhận.
\FormSection{4. Kết quả nghiên cứu:}
Code và hồ sơ local mô tả API/worker/web, Bridge ba placement và cơ chế tenant-aware. Báo cáo phân tích năm nhóm và tổng hợp bằng chứng cùng câu trả lời RQ có điều kiện. Chưa có số đo performance, SUS hoặc kết luận placement tối ưu.
\FormSection{5. Sản phẩm:}
Nguồn ứng dụng, tài liệu kiến trúc/contract/test và nguồn LaTeX/PDF được lưu trong hồ sơ sản phẩm. Tình trạng bàn giao chính thức do nhóm xác nhận.
\FormSection{6. Công bố khoa học hoặc nhận xét của cơ sở áp dụng:}
Chưa có thông tin công bố hoặc nhận xét cơ sở áp dụng được xác nhận trong hồ sơ.
\FormSection{7. Đóng góp và khả năng áp dụng:}
Có thể phục vụ review và học tập kiến trúc đa thuê bao; tác động thực tế, hiệu quả giáo dục và chất lượng dịch vụ chưa được đo.
\FormSection{8. Hiệu quả và phương thức chuyển giao:}
Dự kiến bàn giao mã nguồn/tài liệu sau review của nhóm và người hướng dẫn. Vận hành Internet, provider thật và nghiệm thu nghiên cứu còn cần bằng chứng riêng.
